% Owner: Prisca (Eileen reviews).
\section{Introduction}

LLM-based assistants now review pull requests, explain code and suggest fixes, and developers increasingly
treat their verdict on a piece of code as a security check. Code, however, keeps changing: it is refactored,
documented, logged, split into modules, and sometimes edited by people who would prefer a vulnerability to go
unnoticed. A reviewer that correctly flags a bug today but calls the same bug safe after a routine change gives a
false sense of security that is worse than no review at all.

Prior work on adversarial examples for models of code perturbs programs, for example by renaming identifiers or
inserting dead code, to flip the output of \emph{fine-tuned} vulnerability detectors, typically with white-box or
query access to the model~\cite{yefet2019adversarial,yang2022natural,na2023dip,du2025statement}. We ask a
different question: do \emph{ordinary} changes, of the kind code undergoes anyway, make a general-purpose LLM miss a
bug it caught before? And if they do not, what does?

A central difficulty is ground truth. After an edit, a bug may have been silently fixed, in which case a \safe{}
verdict is correct. We therefore never assume the bug survived: every program the model judges is also executed
by our verifier with attack inputs.

We study four kinds of bug-preserving change, in increasing order of intent: (1) the model's own quality-assurance
(QA) edits to the code; (2) semantics-preserving ``testability tarpit'' transformations that were designed to defeat
program-analysis tools~\cite{alkassar2022tarpits}; (3) embedding the vulnerable handler in a larger, multi-file
application; and (4) deliberate attacks: prompt injection, social pressure, fake test results and unseen helper
functions. We then test a standard prompt-level defense and repeat the key conditions on a commercial model.

\textbf{Contributions.}
\begin{itemize}
\item An execution-verified benchmark of 102 vulnerable/safe JavaScript handlers, plus 1,428 transformed and
      multi-file variants, each checked by the verifier to keep its label.
\item A measurement of four kinds of bug-preserving change on gpt-oss:20b. Natural changes almost never fool it
      (1 \safe{} verdict in 916 still-vulnerable judgments across Parts 1--3), while two deliberate attacks fool it
      more than half the time.
\item Evidence that the two successful attacks differ in kind. The fake end of the code block is a prompt
      injection: a delimited prompt removes it, and a commercial model ignores it. The unseen helper with a
      reassuring name survives the defense and transfers to the commercial model.
\item An open-source harness, verifier, generators and all raw model outputs:
      \url{https://github.com/Prisca-Nova/17739-fooling-the-copilot}.
\end{itemize}
